\documentclass[11pt]{article}
\usepackage[a4paper,margin=2.5cm]{geometry}
\usepackage{amsmath,amssymb,amsthm}
\usepackage[british]{babel}
\usepackage{booktabs}
\usepackage{tabularx}
\usepackage{float}
\usepackage{url}
\newtheoremstyle{result}{\topsep}{\topsep}{\itshape}{}{\bfseries}{.}{ }{\thmnote{#3}}
\theoremstyle{result}
\newtheorem*{result}{Result}
\setlength{\textfloatsep}{12pt plus 2pt minus 4pt}
\newcommand{\E}{\mathbb{E}}
\newcommand{\Q}{\mathbb{Q}}
\renewcommand{\P}{\mathbb{P}}
\newcommand{\HML}{\mathit{HML}}
\newcommand{\FD}{\mathit{FD}}
\newcommand{\fd}{\mathit{fd}}
\newcommand{\rx}{\mathit{rx}}
\newcommand{\Cskew}{C^{\mathrm{skew}}}
\newcommand{\Vsmile}{V^{\mathrm{smile}}}
\newcommand{\Vflat}{V^{\mathrm{flat}}}
\newcommand{\VGK}{V^{\mathrm{GK}}}
\newcommand{\thUB}{\theta_{\mathrm{UB}}}
\newcommand{\sigP}{\hat\sigma^{P}}
\newcommand{\D}{\Delta}

\title{Crash insurance and the G10 carry premium}
\author{Amanjeet Singh}
\date{Draft of 26 September 2026}

\begin{document}
\maketitle

\begin{abstract}
\noindent This paper asks whether the price of crash insurance in G10 currency option smiles accounts for the carry premium and its change between the zero-rate regime and the hiking regime from 2022. Month-end SABR smiles for nine currencies against the dollar, calibrated to dealer quotes from May 2013 to August 2026, price $10\D$ protective options on a three-long, three-short carry portfolio. The point estimate of the skew price of protection per unit of forward discount falls from 0.809 to 0.454, but the pre-registered interval for the difference includes zero. The ratio predicts carry returns in sample (bias-corrected one-sided $p=0.021$) but not out of sample (Clark--West statistic 1.639 against 1.645). Only one of three pre-registered conditions holds, so the regime-shift explanation has partial support only. The unhedged carry premium is itself insignificant ($t=1.47$), so the hedge-cost ratio is not identified: its confidence set is unbounded.
\end{abstract}

\section{Introduction}

A currency carry trade funds positions in high-yield currencies by borrowing in low-yield ones. Its losses are concentrated in bad states. Brunnermeier, Nagel and Pedersen (2008) show that high interest-rate differentials predict negatively skewed currency returns, that carry losses coincide with rises in the VIX and with the unwinding of speculators' positions, and that risk reversals rise after carry losses. If the carry premium compensates for crash risk, the price of insurance against crashes, which foreign-exchange option smiles quote directly, should account for part of it. Burnside, Eichenbaum, Kleshchelski and Rebelo (2011) hedge carry positions with at-the-money options and characterise the peso state as one in which the stochastic discount factor is high. Jurek (2014) forms crash-neutral carry portfolios with out-of-the-money options and finds that crash premia account for at most one third of carry returns. Farhi, Fraiberger, Gabaix, Ranci\`ere and Verdelhan (2015) estimate the share of the carry premium that compensates for disaster risk. The relation has also been unstable: Bussi\`ere, Chinn, Ferrara and Heipertz (2022) find positive Fama-regression coefficients from about 2006 through the decade after the financial crisis, and negative ones again after 2018.

This paper asks whether the price of crash insurance, measured from G10 foreign-exchange option smiles, accounts for the carry premium and for its change between the zero-rate regime (May 2013 to December 2021) and the hiking regime (from January 2022). The regime labels follow the design and denote calendar periods, not rate paths. The ex-ante carry of a three-long, three-short portfolio, its forward-discount spread, averages 18.61 basis points (bp) per month in the first regime and 27.05 in the second (a supplementary, post hoc split). Realised carry returns, compared by regime in a further post hoc analysis, average 7.4 bp per month in the first regime and 37.0 in the second, but the difference (29.6 bp, standard error 23.8) is not significant, and their full-sample mean is itself insignificant ($t=1.47$). The regime comparison therefore concerns the forward-discount spread, and asks whether the option market's price of left-tail protection moved with it.

I fixed the research design on 23 September 2026, before any return, hedge cost or option-implied statistic was computed. Later changes are recorded in the research log with their date and reason, and where a result had been computed under the original rule, it remains reported. Month-end smiles for nine currencies against the dollar are calibrated to dealer-contributed at-the-money, $25\D$ risk-reversal and $25\D$ butterfly quotes from LSEG Workspace. The carry portfolio is formed as in Lustig, Roussanov and Verdelhan (2011), and each leg is protected by a one-month $10\D$ option that pays when the leg loses. The primary estimand (E1), $\varphi$, is the skew premium of that protection (the premium at the smile volatility less that at the at-the-money volatility, same strike) per unit of the portfolio's forward discount. E2 asks whether $\varphi$ predicts carry returns; E3 splits the realised hedge cost into payoff (delta and diffusive), volatility-level and skew terms; E4 repeats E1 and the skew term under the other butterfly reading; E5 measures exposure to VIX-futures roll-down and global FX volatility.

The primary estimate does not support a claim. Under the market-strangle reading with $10\D$ hedges, $\varphi$ averages 0.809 in the zero-rate regime and 0.454 in the hiking regime; the difference of $-0.355$ has a Newey--West standard error of 0.204 and a stationary-bootstrap 95\% interval of $[-0.814,0.014]$, so I do not claim a regime difference. A supplementary split, added after the first results, shows why the point estimate falls: the skew cost of protection is almost unchanged, at 11.24 bp per month in the first regime and 11.19 in the second, while the forward-discount spread rises by 8.44 bp. After the small-sample bias correction, a higher $\varphi$ predicts higher carry returns in sample (one-sided $p=0.021$), the sign a crash-compensation reading implies, but the out-of-sample Clark--West statistic of 1.639 does not reach the 5\% critical value of 1.645. Only one of the three pre-registered conditions for explaining the regime shift is met, and it concerns in-sample predictability rather than the change between regimes, so the explanation has partial support only.

The ex-post evidence is consistent with this reading. Over 2013--2026 the mean realised cost of $10\D$ protection is close to the skew premium paid ex ante ($-11.2$ bp per month), because the payoff and volatility-level terms are indistinguishable from zero. The hedge-cost ratio $\thUB$ is 0.53, against a diffusive null of about 0.10, but its 95\% test-inversion set is unbounded and I claim no value for it. Unhedged carry loads strongly on the VIX roll-down; the hedge lowers the loading by about a fifth, and neither portfolio earns a significant $\alpha$. Among the robustness variants, only a ten-currency ranking in which the dollar can be a leg makes the difference in $\varphi$ significant under both the Newey--West and the bootstrap inference, and no variant rejects out of sample when training starts in 2013. The estimates point to a rise in the forward-discount spread rather than a change in the price of crash insurance.

Section~\ref{sec:framework} sets out what option prices can and cannot identify. It restates the decomposition of the currency premium of Backus, Foresi and Telmer (2001), in the entropy form of Backus, Chernov and Zin (2014), with the covered-interest-parity basis as a wedge, much as Lustig and Verdelhan (2019) do with a general wedge: option prices identify the risk-neutral law of the log exchange-rate change but not how the entropy difference behind the currency premium divides across cumulant orders, so the design measures the price of the carry position's left-tail exposure, not the crash share of discount-factor entropy. I also bound option-implied moments from truncated quotes, in the spirit of the truncation bounds of Jiang and Tian (2005) but under a stated tail condition, and derive the diffusive null of the hedge cost, whose leading-order form is due to Farhi et al.\ (2015), in an exact form with an error bound.

The closest work is Jurek (2014), who studies unhedged carry over 1990--2012 and option-hedged carry over 1999--2012; this paper extends option-hedged evidence for G10 carry through the hiking regime to 2026. Burnside et al.\ (2011) hedge with at-the-money options; the hedged return here follows their construction, but the primary hedge is at $10\D$, because a one-month $25\D$ option is only about 2\% out of the money. Farhi et al.\ (2015) estimate a disaster share and state the leading-order form of the diffusive null (credited in Section~\ref{sec:r6}); this paper estimates prices and exposures, not a disaster share. Chernov, Graveline and Zviadadze (2018) estimate the jump share of the conditional entropy of dollar exchange-rate changes over 1986--2010, which is distinct from the share of the carry premium. Della Corte, Ramadorai and Sarno (2016) use the volatility risk premium as a currency predictor, read the quoted butterfly as a smile strangle and sort currencies on the $10\D$ risk reversal; here the effect of the reading is measured (E4). Fan, Londono and Xiao (2022) show that equity tail risk, measured from one-month out-of-the-money S\&P 500 puts, is priced in currency returns; here the tail price comes from the currencies' own options. Caballero and Doyle (2012) relate carry returns to a VIX futures roll-down strategy, which E5 adapts to month-ends, and argue that exchange-rate options are cheap systemic insurance. Choi and Suh (2022), according to their abstract and introduction, use risk reversals to predict carry crashes and to hedge conditionally; the hedge here is unconditional. Kutuk and van Wijnbergen (2025) apply the hedged-versus-unhedged comparison to USD/TRY only.

The paper adds three pieces of evidence to this literature. The first is a decomposition of the realised cost of G10 crash insurance over 2013--2026 into payoff, volatility-level and skew terms, with a comparison of its ex-ante skew component across the two rate regimes; the design stated this decomposition across both regimes, but the reports give the three-term split for the pooled sample only. The second is a pre-registered test of whether the ex-ante skew price explains the change in the carry premium between the regimes. The third, secondary, is the size of the difference between the two readings of the butterfly quote for the carry estimands; Reiswich and Wystup (2012) set out the market-strangle calibration on which the distinction rests, and Bossens, Ray\'ee, Skantzos and Deelstra (2010, Section 3.3) show that the two readings can differ markedly when the smile is skewed.

Section~\ref{sec:framework} presents the framework, Sections~\ref{sec:data} and~\ref{sec:method} the data and method, Section~\ref{sec:results} the results and Section~\ref{sec:robust} their robustness, and Section~\ref{sec:concl} concludes. Appendix~\ref{app:theory} lists the theory results and Appendix~\ref{app:repro} gives the steps needed to reproduce them.

\section{Framework: what option prices identify}\label{sec:framework}

The results are labelled R1 to R9 as in \texttt{theory/notes.tex}, which contains the proofs; they are stated here with their hypotheses. Appendix~\ref{app:theory} lists their status and sources. Throughout, $\Phi$ is the standard normal distribution function, $\tau$ is the time to expiry and, for forward $F$, strike $K$ and volatility $\sigma$, $d_\pm(K,\sigma)=[\ln(F/K)\pm\sigma^2\tau/2]/(\sigma\sqrt\tau)$; with $s=\sigma\sqrt\tau$ fixed this is written $d_\pm(K)$.

\subsection{Currency premia and discount-factor entropy}

Let $S_t$ be the USD price of one unit of a foreign currency, $s_t=\log S_t$, $f_t$ the log one-period forward, $r_t$ and $r^*_t$ the one-period USD and foreign interest rates, and
\[
x_t=(f_t-s_t)-(r_t-r^*_t)
\]
the covered-interest-parity basis of Du, Tepper and Verdelhan (2018); they quote exchange rates in foreign currency per USD, so $x_t$ is minus their basis. For a positive random variable $M$ with $\E_t|\log M|<\infty$, conditional entropy is $L_t(M)=\log\E_tM-\E_t\log M\ge0$.

\begin{result}[R7 (entropy decomposition)]
Let $M,M^*>0$ be $\mathcal F_{t+1}$-measurable discount factors in $L^2$ that price USD and foreign payoffs in their own currencies, with $\E_tM=e^{-r_t}$, $\E_tM^*=e^{-r^*_t}$, $\E_t|\log M|,\E_t|\log M^*|<\infty$ and $MS_{t+1}/S_t\in L^2$. Suppose $M^*=MS_{t+1}/S_t$, which holds if every foreign payoff in $L^2(\mathcal F_{t+1})$ is traded and priced consistently in both currencies. Write $m=\log M$ and $\D s=s_{t+1}-s_t$. Then:
\begin{enumerate}
\item[(i)] $\E_t[\D s+r^*_t-r_t]=L_t(M)-L_t(M^*)$ and $\E_t[s_{t+1}-f_t]=L_t(M)-L_t(M^*)-x_t$;
\item[(ii)] if the conditional cumulant generating function of $m$ has a Maclaurin series with radius of convergence greater than one, $L_t(M)=\sum_{j\ge2}\kappa_{j,t}/j!$, where $\kappa_{j,t}$ are the conditional cumulants of $m$;
\item[(iii)] for any gross return $R>0$ with $\E_t[MR]=1$ and $\E_t|\log R|<\infty$, $L_t(M)\ge\E_t\log R-r_t$.
\end{enumerate}
\end{result}

The decomposition is due to Backus, Foresi and Telmer (2001): with the exchange rate equal to the ratio of the two kernels (their Proposition 1), the expected excess return on a forward position is the difference between the log of the expectation and the expectation of the log of the kernels (their eq.~(11)), and higher cumulants enter only through that difference (eqs.~(12)--(13)). The name conditional entropy and the bound in (iii) are as in Backus, Chernov and Zin (2014), who relate the bound to Bansal and Lehmann (1997) and Alvarez and Jermann (2005). Lustig and Verdelhan (2019, Section 2.5) restate the decomposition with a wedge between the exchange rate and the ratio of kernels, and Jiang, Krishnamurthy and Lustig (2021) link the value of the dollar to a convenience yield inferred from the Treasury basis, a covered-parity deviation measured with government bond yields. Writing the covered-interest-parity basis into the identity is an elementary step, derived in the notes.

By (i), the expected log excess return on a foreign currency is the difference of two conditional entropies less the basis. By (ii), entropy is a sum of cumulants of log marginal utility, so crash risk enters through cumulants of order three and above. By (iii), applied to the gross returns of fully collateralised hedged and unhedged carry positions, the two positions give two lower bounds on the same $L_t(M)$.

\subsection{Change of numeraire, arbitrage-free prices and spanning}

Options are quoted on currency pairs, five of which have the dollar as the base currency. The change of numeraire connects the two directions of quotation.

\begin{result}[R1 (domestic and foreign risk-neutral measures)]
Let $(\Omega,\mathcal F,(\mathcal F_t)_{t\in[0,T]},\Q^d)$ satisfy the usual conditions, let $B^d,B^f$ be strictly positive money-market accounts with $B^d_0=B^f_0=1$, let $S>0$ be adapted and c\`adl\`ag, and suppose $Z_t=S_tB^f_t/(S_0B^d_t)$ is a true $\Q^d$-martingale. Define $d\Q^f/d\Q^d=Z_T$. Then (a) $\Q^f\sim\Q^d$; (b) $\E^{\Q^f}[X\mid\mathcal F_t]=\E^{\Q^d}[Z_TX\mid\mathcal F_t]/Z_t$ for $\mathcal F_T$-measurable $X\ge0$ or $X\in L^1(\Q^f)$; (c) foreign prices discounted by $B^f$ are $\Q^f$-martingales exactly when their domestic values discounted by $B^d$ are $\Q^d$-martingales; (d) under Garman--Kohlhagen dynamics, $\E^{\Q^d}S_T=F$ and $\E^{\Q^f}(1/S_T)=1/F$.
\end{result}

This is the change of numeraire of Geman, El Karoui and Rochet (1995), proved in the notes for completeness. If $Z$ is only a strict local martingale, $Z_T$ does not define an equivalent measure; Carr, Fisher and Ruf (2014) construct a foreign measure in that case that is not equivalent to $\Q^d$. In the moment calculations of Section~\ref{sec:e2}, prices on a pair whose quote currency is not the dollar are converted to the USD forward measure with density $S_T/F$, where $S$ is that pair's price of one dollar.

\begin{result}[R2 (arbitrage-free call prices)]
Let $c:[0,\infty)\to\mathbb R$ and $F>0$. There is a random variable $X\ge0$ with $\E X=F$ and $c(K)=\E(X-K)^+$ for all $K\ge0$ if and only if $c$ is convex, $c(0)=F$, $c'_+(0)\ge-1$ and $c(K)\to0$ as $K\to\infty$. The law of $X$ is then unique, with distribution function $1+c'_+$ on $[0,\infty)$.
\end{result}

The link between the second strike derivative of call prices and state prices is due to Breeden and Litzenberger (1978); the characterisation is a known result, proved by the same argument in F\"ollmer and Schied (2004, Lemma 7.23), and is restated in the notes. On a strike grid the necessary conditions reduce to slopes in $[-1,0]$ that are nondecreasing, which is the static-arbitrage test used below, together with the density condition on total implied variance of Gatheral and Jacquier (2014). An arbitrage-free smile at one maturity therefore identifies the whole risk-neutral law of the terminal exchange rate. The quotes, however, pin the smile down only at the at-the-money, $25\D$ and $10\D$ strikes.

\begin{result}[R3 (spanning)]
Let $f:(0,\infty)\to\mathbb R$ have a locally absolutely continuous derivative, let $X>0$ almost surely with $\E X=F<\infty$, and write $P(K)=\E(K-X)^+$ and $C(K)=\E(X-K)^+$. If $\int_0^F|f''|P\,dK+\int_F^\infty|f''|C\,dK<\infty$, then $f(X)\in L^1$ and
\[
\E f(X)=f(F)+\int_0^F f''(K)P(K)\,dK+\int_F^\infty f''(K)C(K)\,dK.
\]
\end{result}

The spanning formula is that of Carr and Madan (2001); the notes write out the standard Taylor-remainder proof with explicit integrability conditions under which expectations can be taken term by term. The volatility, cubic and quartic contracts of Bakshi, Kapadia and Madan (2003) are of this type. They expand around the spot price; the contracts used here are centred at the forward.

\subsection{Partial identification from truncated quotes}

Quotes observe $P$ and $C$ only between the $10\D$ strikes $K_{\min}<F<K_{\max}$, so the tails of the spanning integrals must be bounded.

\begin{result}[R4 (bounds on the unobserved tails)]
Let $G(u)=\Q(X\le u)$. (a) If $G(u)/u^\gamma$ is nondecreasing on $(0,K_{\min}]$ for some $\gamma>0$, then $P(K)\le P(K_{\min})(K/K_{\min})^{\gamma+1}$ for $0<K\le K_{\min}$. (b) If $u^\eta\Q(X>u)$ is nonincreasing on $[K_{\max},\infty)$ for some $\eta>1$, then $C(K)\le C(K_{\max})(K/K_{\max})^{1-\eta}$ for $K\ge K_{\max}$. Hence, for any measurable weight $w\ge0$,
\[
\int_0^{K_{\min}}w(K)P(K)\,dK\le\frac{P(K_{\min})}{K_{\min}^{\gamma+1}}\int_0^{K_{\min}}w(K)K^{\gamma+1}dK,
\]
and similarly in the upper tail. In particular,
\[
\int_0^{K_{\min}}\frac{P(K)}{K^2}\,dK\le\frac{P(K_{\min})}{\gamma K_{\min}},\qquad \int_{K_{\max}}^\infty\frac{C(K)}{K^2}\,dK\le\frac{C(K_{\max})}{\eta K_{\max}}.
\]
\end{result}

The bounds use only observed wing prices and assumed tail exponents. Sign-changing weights are split into positive and negative parts. Because $\gamma$ and $\eta$ are not observed, risk-neutral variance and skewness are partially identified: they are reported as intervals over stated exponents, never as point values. Jiang and Tian (2005, Proposition 2) bound the truncation error of model-free implied variance beyond the quoted strikes; the bounds here rest on a different, stated tail condition and extend to the higher moment contracts. The lemma, the bounds and their proofs are my own and are in the notes.

\subsection{The diffusive null of the hedge cost}\label{sec:r6}

As Jurek (2014) notes, buying protection costs part of the carry premium even when no crash risk is priced, because an unlevered option hedge also gives up part of the diffusive premium. For a forward $G$ and the put's strike $K$, write $d_\pm(G)=[\ln(G/K)\pm s^2/2]/s$.

\begin{result}[R6 (diffusive null of the hedge cost)]
Suppose that under $\Q$ the terminal rate $S_\tau$ is lognormal with mean $F$ and log-variance $s^2$, and that under $\P$ it is lognormal with the same log-variance and mean $Fe^{\lambda\tau}$, $\lambda\ne0$, with no jumps and no volatility premium. A unit long forward has $\P$-expected payoff $\mu_U=F(e^{\lambda\tau}-1)$; adding a put struck at $K$, bought at its $\Q$-value with the premium financed to $\tau$, gives expected payoff $\mu_H$. Then $\theta_0=1-\mu_H/\mu_U=\Phi(-d_+(\xi))$ for some $\xi$ between $F$ and $Fe^{\lambda\tau}$, and
\[
\bigl|\theta_0-\Phi(-d_1)\bigr|\le\frac{1}{\sqrt{2\pi}}\,\frac{|\lambda|\sqrt\tau}{\sigma},\qquad d_1=d_+(F).
\]
\end{result}

Farhi et al.\ (2015, Section 5.1) state the leading-order form: a put of delta $\D_P$ leaves a fraction $1+\D_P$ of the Gaussian risk premium. Jurek (2014) reads the hedged-to-unhedged comparison as an upper bound on the crash share. The mean-value form and the error bound are my own short derivation. The notes treat a long forward hedged with a put; a short forward hedged with a call follows by the same argument, with the call's forward delta in place of the put's. With no crash premium at all, a put of forward delta $-\Phi(-d_1)$ surrenders approximately that fraction of the carry premium: about 0.25 for a $25\D$ put and 0.10 for a $10\D$ put. This is the null $\theta_0$ against which the realised ratio $\thUB$ of Section~\ref{sec:method} is read.

\subsection{What the design measures}

By R7(i), the carry premium is an entropy difference less the basis, and crash risk enters it through higher cumulants. By R1, R2 and R4, option prices on $S$ reveal the risk-neutral law of $\D s=m^*-m$ at the quoted maturity, only partially beyond the $10\D$ strikes, but neither the physical law nor the separate cumulants of $m$ and $m^*$; they constrain, but do not identify, how $L_t(M)-L_t(M^*)$ divides across cumulant orders. The design accordingly measures the price of the carry position's left-tail exposure, through the premia of options that pay when the legs lose, and does not measure the crash share of discount-factor entropy. A premium that remains after hedging can reflect diffusive risk (R6); options that are cheap as systemic insurance, which is how Caballero and Doyle (2012) explain why option-hedged carry retains high returns that systemic risk does not account for; or the pricing of a peso state (Burnside et al., 2011). The design does not separate these.

\section{Data}\label{sec:data}

\subsection{Instruments and contributors}

I obtained the option and money-market data from LSEG Workspace under a student licence provided by my university (Section~\ref{sec:licence}). The main retrieval, of 23 September 2026, requested daily history from 1 January 1995. For each of the nine currencies (EUR, GBP, AUD, NZD, JPY, CHF, CAD, NOK and SEK, all against USD) the panel contains spot and one-month forward points (composite, bid and ask); at-the-money volatility, $25\D$ and $10\D$ risk reversals and $25\D$ and $10\D$ butterflies at one and three months, from the composite and from the Fenics contributor (bid, ask and mid); and one-month deposit rates, with OIS where an instrument exists (none exists for EUR or SEK). For USD it contains one-month Fed Funds OIS, SOFR OIS and a deposit rate. Three-month forward points and rates were retrieved on 24 September 2026 for the three-month variant, and continuation series of the first two VIX futures complete the panel. Named-broker contributors are short, sparse or mid-only; one is used only as a check. WM/Reuters fixings, the S\&P 500 index and the Cboe VIX index are not licensed. Mid quotes are the average of bid and ask.

The provider does not document the time of day of daily history. Matching daily values against intraday bars in winter and summer 2026 places the composite value at about 21:18 UTC throughout the year and the Fenics value at about 17:16 London time. Option payoffs use end-of-day spot on the expiry date, because the 10:00 New York cut is not available; a robustness variant uses the previous day's spot.

\subsection{Samples and the month-end rule}

Sampling is at New York month-ends: the last business day of the calendar month under the Federal Reserve holiday calendar (Federal Reserve Board, K.8), with a Saturday holiday not moved to the Friday. A quote is observed if it is two-sided on the month-end and substituted if the last two-sided quote lies within the five preceding business days. Otherwise it is missing and the currency is excluded that month; a month is used only if at least twice the number of legs remain. Substitutions and exclusions are counted in the private audit; the counts for the one-month sample are not reported here.
The primary sample uses composite quotes. It starts at the first month-end at which all 45 composite one-month smile series (at-the-money, $25\D$ and $10\D$ risk reversals and butterflies for nine currencies) and the nine forward-point series are available under this rule, and runs from May 2013 to August 2026: 160 month-ends, of which 104 fall in the zero-rate regime (May 2013 to December 2021) and 56 in the hiking regime (from January 2022). The return window of the last month-end ends after the last spot observation and is excluded, so the return sample has 159 month-ends, May 2013 to July 2026, with returns realised from June 2013 to August 2026.

The extended sample uses Fenics quotes before the primary start, with two long and two short legs, from January 2007 to April 2013; 72 of these month-ends have the required four currencies with the calibration set (at-the-money, $25\D$ risk reversal and butterfly) and forward points. Fenics quotes before 2010 update infrequently, butterflies in particular, so results from this sample are secondary evidence. The design also defines a long at-the-money sample, at-the-money-hedged carry as in Burnside et al.\ (2011) from the start of the at-the-money and forward series; the reports give no result from it, so none is reported here.

\subsection{Quote conventions and audit outcomes}

Quote conventions follow Reiswich and Wystup (2012), who report the pair conventions of Clark (2011): spot delta up to one year; premium-adjusted delta for the dollar-base pairs (USDJPY, USDCHF, USDCAD, USDNOK and USDSEK) and pips delta otherwise; at-the-money as the delta-neutral straddle, with strike $F\exp(\pm\sigma^2\tau/2)$ for pips and premium-adjusted delta respectively; and the risk reversal as call minus put volatility on the quoted pair. For dollar-base pairs, spot and forward are inverted so that $X^j_t$ is the USD price of one unit of currency $j$; the outright forward is spot plus forward points times the pip factor.

The data audit checked coverage, two-sidedness, staleness, the underlying pair, quote semantics, conventions, the splice between contributors and plausibility (implausible quotes are flagged, not removed). Provider metadata give every volatility instrument, including those for NOK and SEK, as the currency against USD, and a realised-volatility test agrees, so no triangulation through EURUSD is needed. The forward-point scaling factors equal the configured pip factors. Risk reversals are call minus put on the quoted pair, and the $10\D$ quotes follow the $25\D$ convention. Instrument metadata name only the quote type and the delta, and LSEG's documentation of its own surfaces does not define the contributed quotes (LSEG Developer Portal), so delta and premium conventions follow market convention as above. Fenics and composite skew costs agree leg by leg, with correlations of 0.995 before the primary sample and 0.997 within it. A second retrieval, needed to check for quote revisions, has not yet been made.

\subsection{The butterfly reading}\label{sec:bf}

The quoted $25\D$ butterfly $BF$ can be read in two ways. Under the market-strangle reading, a strangle struck at the $25\D$ strikes of the flat volatility $\sigma_{ATM}+BF$ has the same premium under the smile as at that flat volatility. Under the smile-strangle reading, $BF$ is the average of the smile's volatilities at its own $25\D$ strikes less $\sigma_{ATM}$. A three-parameter smile fits either reading exactly, so the quotes cannot decide between them. The design sets the reading by provider documentation first, then by the interdealer market-strangle convention (Reiswich and Wystup, 2012), with the fit to the held-out $10\D$ quotes as supporting evidence only. Provider documentation is silent, so the market-strangle reading is primary. The supporting diagnostic points the other way: smiles calibrated under the smile-strangle reading predict the held-out $10\D$ butterflies better in most month-ends and in every currency, with small absolute differences. This is recorded as evidence against the primary reading, and E4 (smile-strangle reading) is reported alongside the primary results.

\subsection{VIX futures}

Individual VIX futures contracts are taken from the Cboe Futures Exchange historical settlement files, one per contract, for contracts from January 2007 to December 2026. Before 26 March 2007 the archive quotes contracts at ten times the index level; the change of scale is detected in the files and removed. After this, Cboe second-month settlements equal the LSEG second-month continuation at every month-end (223 of 223). The Cboe files are stored privately because their redistribution terms are not established.

\subsection{Licence}\label{sec:licence}

I obtained the LSEG data under an LSEG Workspace student licence provided by my university. The licence covers only me: it permits individual study and research, and it does not permit redistribution of the data or placing them in a public repository. Neither this paper nor the repository therefore contains LSEG data. They publish code, methods and aggregate results: means, standard errors, test statistics, intervals and coverage facts such as sample windows and counts. Quote values, month-level series and calibrated parameters are held privately, and no figure of LSEG-derived time series is shown. Anyone else who wants to rerun the acquisition, and hence the estimation, needs their own LSEG Workspace licence.

\section{Method}\label{sec:method}

\subsection{Smile calibration}

For each currency, month-end and butterfly reading, a SABR smile with $\beta=1$ (Hagan, Kumar, Lesniewski and Woodward, 2002) is calibrated to the at-the-money volatility, the $25\D$ risk reversal and the $25\D$ butterfly. Under the primary reading the three equations are: the smile's volatility at the delta-neutral-straddle strike equals $\sigma_{ATM}$; the difference of the smile's volatilities at its own $25\D$ call and put strikes equals the risk reversal; and the market strangle has the same premium under the smile as under the flat volatility $\sigma_{ATM}+BF$ (Reiswich and Wystup, 2012). The system is solved by Levenberg--Marquardt in $(\ln\alpha,\operatorname{atanh}\rho,\ln\nu)$, starting from the previous month's solution with three fixed starting points as fall-backs. A calibration has converged when the largest absolute residual is below $10^{-8}$ volatility units. The $10\D$ quotes are held out as a fit test. Strikes are found from deltas in each pair's own convention.

\begin{result}[R5 (delta-to-strike inversion)]
(a) For pips spot delta $\D=\omega D_b\Phi(\omega d_+)$ of a call ($\omega=1$) or put ($\omega=-1$) at a flat volatility, with $D_b$ the base-currency discount factor, $K=F\exp\bigl(\sigma^2\tau/2-\omega\sigma\sqrt\tau\,\Phi^{-1}(|\D|/D_b)\bigr)$. (b) The premium-adjusted call delta $h(K)=D\,(K/F)\,\Phi(d_-(K))$, with $D$ the base-currency discount factor for spot delta and $D=1$ for forward delta, tends to zero at both ends, is strictly increasing up to a unique maximiser $K^*$ and strictly decreasing after it; each $\D\in(0,h(K^*))$ is attained at exactly two strikes, and the convention takes the one in $[K^*,\infty)$. (c) For a continuously differentiable smile $\sigma(K)$, if $|\partial\sigma/\partial k|\sqrt\tau\,|d_-(K,\sigma(K))|<1$ on an interval, with $k=\ln(K/F)$, the pips delta $\omega D\Phi(\omega d_+(K,\sigma(K)))$ is strictly decreasing there, so each delta attained on the interval is attained at exactly one strike.
\end{result}

Reiswich and Wystup (2012) give the delta definitions, the non-monotonicity of the premium-adjusted call delta and the right-branch convention. J\"ackel (2020, Section 2, eq.~(20)) gives the single maximum in (b) through the inverse Mills ratio, the argument the notes write out. Reiswich (2010, Section 3.8) shows that a smile delta can fail to be monotone; the sufficient condition in (c) is my own. Across the calibrated one-month smiles of the primary sample, the condition in (c) holds over $\pm4$ at-the-money standard deviations at every month-end, and premium-adjusted deltas are monotone on the conventional branch everywhere.

\begin{result}[R8 (local well-posedness of the calibration)]
For $\beta=1$ the Hagan implied volatility is infinitely differentiable in $(K,\alpha,\rho,\nu)$ on $\{K>0\}\times(0,\infty)\times(-1,1)\times[0,\infty)$. Let $\Psi(\theta,q)=0$ be the three calibration equations in $\theta=(\alpha,\rho,\nu)$ for quotes $q=(\sigma_{ATM},RR,BF)$, with a solution $(\theta^*,q^*)$. Suppose that at the solution (i) the at-the-money strike equation has a non-zero strike derivative; (ii) the smile delta has a non-zero strike derivative at the $25\D$ strikes; (iii) for premium-adjusted delta, the $25\D$ call strike at the flat volatility $\sigma_{ATM}+BF$ lies strictly right of the maximiser $K^*$; and (iv) $D_\theta\Psi$ is non-singular. Then there are neighbourhoods $U\ni q^*$ and $W\ni\theta^*$ and a unique $C^1$ map $\theta:U\to W$ with $\Psi(\theta(q),q)=0$.
\end{result}

The proof applies the implicit function theorem first to the strikes and then to the system, following the argument Reiswich (2010, Theorems 2 and 3) uses for the market-strangle calibration of a simpler smile. Le Floc'h and Kennedy (2014, Section 4, eqs.~(30)--(31)) invert the leading-order expansion of the Hagan formula (Hagan et al., 2002, eq.~(3.1a)) in closed form, recovering $(\alpha,\rho,\nu)$ from the at-the-money level, skew and curvature, which is an approximate, global counterpart of the local and exact statement here. Because the result is local, the solver restarts from several points; condition (iv) is checked at every calibration through the solver's Jacobian, whose condition numbers confirm non-singularity everywhere. Under the smile-strangle reading the result holds without (iii).

Every smile is checked for static arbitrage (R2) on log-moneyness grids of $\pm4$ and $\pm10$ at-the-money standard deviations, by discrete convexity of call prices and by the density condition of Gatheral and Jacquier (2014). Non-convergence and arbitrage flags are counted, not dropped. In the primary sample every month-end smile converged under both readings, with residuals below $10^{-13}$, and every smile passed both checks.

\subsection{Portfolios and protective options}

At each month-end $t$ the available currencies are ranked by the forward discount $\fd^j_t=s^j_t-f^j_t=\log(X^j_t/F^j_t)$, where $F^j_t$ is the one-month outright forward. Lustig, Roussanov and Verdelhan (2011) sort currencies into six portfolios on the forward discount at each month-end and take the highest minus the lowest; with nine currencies the analogue is long the top $n=3$ and short the bottom $n=3$, each with $1/n$ USD of forward notional ($n=2$ in the extended sample). With arithmetic excess returns $\rx^j_{t+1}=X^j_{t+1}/F^j_t-1$ per USD of forward notional,
\[
\HML^U_{t+1}=\frac1n\sum_{\text{long}}\rx^j_{t+1}-\frac1n\sum_{\text{short}}\rx^j_{t+1},\qquad
\FD_t=\frac1n\Bigl(\sum_{\text{long}}\fd^j_t-\sum_{\text{short}}\fd^j_t\Bigr)\ge0.
\]
Returns are evaluated at the spot quote on the option expiry date, the spot rate for value on the forward's delivery date. Dates follow the New York business calendar, and volatility time is trade to expiry in days over 365.

Each leg buys a one-month option that pays when the leg loses: a put on the currency for a long leg, which is a USD call for dollar-base pairs, and a call on the currency for a short leg. The hedge follows the construction of Burnside et al.\ (2011, Section 1.2): option notional matches the forward notional, and premia are compounded forward, here to the delivery date. Premia are therefore in USD per USD of forward notional at delivery, which for both quotation types equals the undiscounted premium on the quoted pair divided by the pair forward. The primary strike is the $10\D$ strike on the calibrated smile in the pair's own delta convention (R5); $25\D$ and at-the-money hedges are comparisons. Pricing is Garman and Kohlhagen (1983) in forward form. The quote-currency discount factor comes from Fed Funds OIS for USD and one-month deposit rates otherwise, and the base-currency factor is implied by the forward, $D_b=FD_q/S$, so no separate covered-interest-parity assumption is made. With $\Pi_{t+1}$ the option payoff per USD of notional and $\Vsmile_t$ its premium at the smile volatility,
\[
\HML^H_{t+1}=\HML^U_{t+1}+\frac1n\sum_{\text{legs}}\bigl(\Pi_{t+1}-\Vsmile_t\bigr).
\]
Estimands use mid prices. An implementable version trades forwards at bid and ask and buys options at the smile volatility plus $k\in\{1,1.5,2\}$ times the at-the-money half-spread, converted to premium by vega, because package quotes do not give a single-option ask.

\subsection{Estimands}

The primary estimand, E1, is measured ex ante. For each leg, $\Vsmile_t$ is the premium of the $10\D$ protective option at the calibrated smile volatility and $\Vflat_t$ the premium at the same strike with the at-the-money volatility. Then
\[
\Cskew_t=\frac1n\sum_{\text{legs}}\bigl(\Vsmile_t-\Vflat_t\bigr),\qquad \varphi_t=\Cskew_t/\FD_t,
\]
the skew price of protection per unit of forward discount. E1 reports the means of $\varphi_t$ in the zero-rate and hiking regimes and their difference. The $25\D$ version is a comparison. The at-the-money comparison named in the design is identically zero, because the calibrated smile passes through the at-the-money volatility, and is not reported.

E2 asks whether the skew price predicts carry returns, through the regression $\HML^U_{t+1}=a+b\,z_t+e_{t+1}$ with $z_t=\varphi_t$ (primary) or the oriented $10\D$ risk reversal: the quoted volatility of the protective option minus that of the opposite option, averaged over legs. The out-of-sample test uses an expanding window, with the first forecast, made at the end of December 2016, for January 2017 and training that starts with the extended sample. Option-implied variance and skewness are secondary predictors, reported as R4 intervals.

E3 is measured ex post. The realised excess payoff of each option splits into three terms, all carried to delivery:
\[
\Pi-\Vsmile=\underbrace{\bigl[\Pi-\VGK(\sigP)\bigr]}_{\text{(i) payoff}}\;\underbrace{-\bigl[\Vflat-\VGK(\sigP)\bigr]}_{\text{(ii) volatility level}}\;\underbrace{-\bigl[\Vsmile-\Vflat\bigr]}_{\text{(iii) skew}},
\]
where $\VGK(\sigP)$ is the Garman--Kohlhagen premium at the same strike with $\sigP$, the annualised realised volatility of daily log spot changes over the previous 21 business days. Term (i) carries the delta and diffusive part; term (iii) is known at $t$. Averaged over legs with weight $1/n$, the three terms add up to $\HML^H-\HML^U$. The hedge-cost ratio is
\[
\thUB=1-\hat\mu_H/\hat\mu_U,
\]
with $\hat\mu_H$ and $\hat\mu_U$ the sample means of $\HML^H$ and $\HML^U$; it can take any real value. Jurek (2014) reads such a ratio as an upper bound on the crash share, but an unlevered hedge also gives up part of the diffusive premium, so $\thUB$ is compared with $\theta_0$ of R6 and is not described as a crash share. Here $\theta_0$ is the equal-weight average over legs of the absolute forward delta of the hedging options; for the portfolio, R6 weights each leg by its expected forward payoff, so the equal weights are an approximation.

E4 recomputes E1 and the skew term of E3 under the smile-strangle reading of the butterfly quote.

E5 measures exposure to systemic risk: $\HML^U$ and $\HML^H$ are regressed on two factors, both measured over each carry return window from the month-end to the one-month option expiry. The VIX roll-down is minus the percentage change of the second-month VIX future's settlement price, holding the same contract. It is a month-end analogue of the strategy of Caballero and Doyle (2012), who short the VIX future whose expiry matches the one-month forward's maturity and hold it to expiry. The global FX-volatility level follows Menkhoff, Sarno, Schmeling and Schrimpf (2012, eq.~4): the average over the window's days of the cross-sectional mean absolute daily log spot change of the nine currencies. Its innovation is the residual of an AR(1) fitted to that series. E5 reports the exposure-adjusted ratio $1-\alpha_H/\alpha_U$.

The design also names a descriptive account of smiles, $\varphi_t$ and returns around 5 August 2024. It consists of daily and month-level values and is not reported.

\subsection{Inference}

Standard errors are those of Newey and West (1987) with a Bartlett kernel, and the bandwidth follows the automatic procedure of Newey and West (1994) without prewhitening: a pilot truncation of $\lfloor4(T/100)^{2/9}\rfloor$, then a bandwidth of $\lfloor\hat\gamma T^{1/3}\rfloor$. The E1 difference is the coefficient on a regime dummy. Its bootstrap interval uses the stationary bootstrap of Politis and Romano (1994), resampling each regime separately with its own expected block length chosen by the rule of Politis and White (2004) as corrected by Patton, Politis and White (2009), with 9,999 draws and percentile intervals. The tuning caps of the block-length rule are taken from the arch package (Sheppard, 2026), which I read to check formulas and constants; no code is copied.

For the predictive regressions, a residual bootstrap under the null of no predictability, with 9,999 draws, draws the pairs $(\hat u_k,\hat v_{k+1})$ independently, where $\hat u$ are residuals of the unrestricted regression and $\hat v$ are AR(1) residuals of the predictor. The slope is corrected by the bootstrap bias estimate, which is checked against the first-order approximation $-(\sigma_{uv}/\sigma_v^2)(1+3\rho)/T$ of Stambaugh (1999, eq.~18). The out-of-sample test is the one-sided MSPE-adjusted statistic of Clark and West (2007) against the historical mean, with the least-squares standard error they recommend for one-step forecasts, since the design names none; the Newey--West version is reported alongside.

The 95\% set for $\thUB$ is obtained, as the design specifies, by test inversion: $\{\vartheta:|\hat\mu_H-(1-\vartheta)\hat\mu_U|\le1.96\,\widehat{se}_{HAC}(\vartheta)\}$, with $\widehat{se}_{HAC}(\vartheta)$ the Newey--West standard error of $\hat\mu_H-(1-\vartheta)\hat\mu_U$. This is the construction of Fieller (1954, Sections 2--4) for a ratio of means, with a Newey--West variance in place of his independent variance estimate: the set is a bounded interval only when the denominator $\hat\mu_U$ is significantly different from zero, and is otherwise the complement of an interval or the whole line. Unbounded sets are reported as unbounded. The design names delta-method intervals as secondary, and none is reported. Because E1 is the single primary estimand, no multiple-testing adjustment is made.

\begin{result}[R9 (inference; cited)]
(a) If $(X_t)$ is strictly stationary with $\E|X_t|^{2+\delta}<\infty$ for some $\delta>0$, strong-mixing coefficients with $\sum_n\alpha(n)^{\delta/(2+\delta)}<\infty$ and a positive long-run variance, its sample mean is asymptotically normal (Ibragimov, 1962, Theorem 1.7). (b) The stationary bootstrap is valid under the moment and mixing conditions of Theorem 2 of the technical-report version of Politis and Romano (1994); Gon\c{c}alves and de Jong (2003) weaken them. (c) Confidence procedures for a ratio that are bounded with probability one have zero worst-case coverage (Gleser and Hwang, 1987; Dufour, 1997), hence test inversion.
\end{result}

These conditions are assumptions, not checked properties. The persistence of $\varphi$ and the heavy right tail of its zero-rate distribution bear directly on them, and each regime mean averages a short, dependent stretch of calendar time; the post hoc analyses of Section~\ref{sec:e1} address this only in part.

\subsection{Interpretation rules}

Section 9 of the design fixes how results are read. A regime difference in E1 is claimed only if its 95\% interval excludes zero, and its size is read relative to $\FD_t$. The design reports a Newey--West and a stationary-bootstrap interval without saying which governs; I treat a difference as significant only when both exclude zero. This reading is mine, not the design's; I adopted it when interpreting the robustness grid, after the primary results, and the research log records it. It does not affect the primary E1 result, where both intervals include zero, but it decides the dollar-carry variant and the three three-month phases, in which only the bootstrap interval excludes zero (Section~\ref{sec:robust}). The statement that crash-insurance pricing explains the regime shift requires all three of: the E1 difference, $b>0$ in E2 after the Stambaugh correction, and an out-of-sample Clark--West rejection at 5\%; otherwise the result is reported as partial or no support. $\thUB$ is not described as a crash share, and a skew term whose interval includes zero is reported as no evidence of priced crash risk beyond the diffusive and volatility-level terms. Null results are reported with the same prominence as positive ones.

\subsection{Validation}

Pricing and delta conventions are checked against analytic identities (put--call parity, the delta-neutral straddle and delta-to-strike round trips), and smile calibration by synthetic recovery under both readings and both delta types. QuantLib 1.43, an independent open-source implementation, agrees with the project's premia, deltas in three conventions, strike inversion, delta-neutral strikes and Hagan volatilities to between $10^{-9}$ and $10^{-13}$ on the test grid. Before the E1 results were recorded, a separate review of the estimand and inference code rebuilt every leg in QuantLib, which agreed to machine precision, and found that the standard errors matched the R package sandwich (Zeileis and contributors, 2026) exactly. The base variant of the robustness grid reproduces the E1 regime difference and the means of $\HML^U$, $\HML^H$ and the skew term to machine precision.

\subsection{Deviations, corrections and items not reported}\label{sec:dev}

The research log records every departure from the design with its date. Before any inference, the Newey--West lag, first set to the pilot truncation of Newey and West (1994), was replaced by the full automatic procedure, and the extended sample was clarified to require the calibration set only. Before the E1 results were recorded, a review led to New York calendar dates, provider days-to-maturity instead of a 30-day fallback, exclusion of currencies rather than whole months when inputs are missing, and a leg-by-leg splice test; each had a negligible effect on the primary estimates. Deposit rates are used for every non-USD currency, although the design allows OIS where available; the effect on the base-currency discount factor, the only channel through which rates enter E1, is bounded at about 0.1 per cent. Before the E2 and E3 results were recorded, two errors were corrected: the first E2 bootstrap resampled blocks of pairs, which reintroduced the sample's predictive relation into the null and inflated the bias about tenfold relative to Stambaugh's approximation; and the leg weights of $\theta_0$ summed to two. After the first E2 and E3 results, and before the E5 results were recorded, return windows ending after the last spot observation were excluded; the E2 and E3 estimates changed only in the third significant figure, the Clark--West statistic from 1.640 to 1.639. The supplementary E1 analyses are post hoc. The open rules of the robustness grid were fixed after the primary results were recorded but before any variant was estimated, and the choices for the three-month tenor, vanna--volga smiles and moment intervals before those estimates were computed.

Several items named in the design have no result in the reports and are not reported: the long at-the-money sample, the sign and magnitude check of the carry portfolio against published currency portfolios, the SOFR OIS check, the delta-method intervals, the outcome of the check against the named-broker contributor, and the outcome of the closed-form checks of the moment code.

\section{Results}\label{sec:results}

All results are aggregate statistics over the primary sample unless stated. Means and E3 terms are in basis points (bp) of notional per month; the units of regression coefficients are given with each table.

\subsection{E1 and E4: the skew price of protection across regimes}\label{sec:e1}

\begin{table}[t]
\centering\footnotesize
\setlength{\tabcolsep}{4pt}
\caption{E1 and E4: regime means and differences}\label{tab:e1}
\begin{tabular}{llccccc}
\toprule
\multicolumn{7}{l}{\emph{Panel A: pre-registered estimand $\varphi$ and comparisons}}\\
Reading, hedge & Quantity & Zero-rate & Hiking & Difference & HAC s.e. & Bootstrap 95\% \\
\midrule
Market, $10\D$ (E1) & $\varphi$ & 0.809 (0.185) & 0.454 (0.066) & $-0.355$ & 0.204 & $[-0.814,\,0.014]$ \\
Market, $25\D$ & $\varphi$ & 0.613 (0.152) & 0.341 (0.048) & $-0.272$ & 0.166 & $[-0.642,\,0.027]$ \\
Smile, $10\D$ (E4) & $\varphi$ & 0.781 (0.178) & 0.437 (0.062) & $-0.345$ & 0.196 & $[-0.787,\,0.010]$ \\
Smile, $25\D$ & $\varphi$ & 0.595 (0.148) & 0.330 (0.046) & $-0.265$ & 0.161 & $[-0.624,\,0.024]$ \\
\midrule
\multicolumn{7}{l}{\emph{Panel B: supplementary (post hoc) split of the E1 ratio}}\\
Market, $10\D$ & $\Cskew$ & 11.24 (0.82) & 11.19 (0.83) & $-0.05$ & 1.30 & $[-2.44,\,2.20]$ \\
 & $\FD$ & 18.61 (1.77) & 27.05 (1.41) & 8.44 & 2.30 & $[3.49,\,13.65]$ \\
\bottomrule
\end{tabular}
\par\medskip
\begin{minipage}{0.97\linewidth}\footnotesize
Market: market-strangle reading; smile: smile-strangle reading. Primary sample, May 2013 to August 2026: 104 zero-rate and 56 hiking month-ends. $\varphi$ is dimensionless; $\Cskew$ and $\FD$ in bp per month. Newey--West standard errors in parentheses. The interval is the stationary-bootstrap 95\% percentile interval for the difference (hiking minus zero-rate), 9,999 draws. Panel B was added after the first results and does not replace the pre-registered estimand.
\end{minipage}
\end{table}

Table~\ref{tab:e1} reports E1 and E4. Under the market-strangle reading with $10\D$ hedges, the point estimate of $\varphi$ falls from 0.809 to 0.454, but both intervals for the difference include zero, under both readings and at both hedge deltas, so the regime difference is not claimed. The design reads its size relative to $\FD_t$. A supplementary (post hoc) split, Panel B, separates the two components of the ratio: there is no evidence that the skew cost of protection changed (difference $-0.05$ bp, interval $[-2.44,2.20]$), whereas the forward-discount spread is higher in the hiking regime (difference 8.44 bp, interval $[3.49,13.65]$). The point estimate of the fall in $\varphi$ therefore reflects the forward-discount spread, which was smallest in 2020--2021: excluding those years, the difference is $-0.03$ (post hoc, below).

Under the smile-strangle reading (E4) the difference is $-0.345$ (standard error 0.196, interval $[-0.787,0.010]$); the reading moves $\varphi$ by less than 4 per cent and the skew term of E3 from $-11.2$ to $-10.8$ bp per month (Table~\ref{tab:rob}), so the bias from misreading the butterfly convention is small and changes no conclusion.

I added several supplementary analyses after the first E1 results, and they are reported as post hoc. They were prompted by two features of the series: $\varphi$ is highly persistent and its zero-rate distribution has a heavy right tail from months with very small forward discounts; they do not replace the pre-registered estimand. The first-order autocorrelations of $\varphi$ are 0.85 and 0.76 in the two regimes, leaving roughly 8 effective observations per regime. With so few, the automatic Newey--West bandwidth is capped well below what the persistence would require, so the pre-registered intervals are likely too narrow. An AR(1) plug-in standard error of the difference is 0.27 ($t=-1.31$). Excluding 2020--2021, the zero-rate period with the smallest forward discounts, the difference is $-0.03$. The median of $\varphi$ falls from 0.53 to 0.38 and the 10\% trimmed mean from 0.63 to 0.41; the mean of $\log\varphi$ falls from $-0.46$ to $-0.89$, and the ratio of regime means, mean $\Cskew$ over mean FD, from 0.60 to 0.41. The Newey--West standard error of the difference rises with the lag, from 0.079 at lag 0 to 0.204 at lag 9 and 0.236 at lag 20, and the two-sided bootstrap $p$-value is 0.064. Dropping the largest 1, 3, 5 and 10 zero-rate observations gives differences of $-0.32$, $-0.28$, $-0.23$ and $-0.14$. The excluded 2020--2021 period contains 24 of the 104 zero-rate month-ends.

The extended sample gives secondary evidence only, because Fenics butterflies before 2010 are often stale. Over January 2007 to April 2013 (72 month-ends with at least four currencies, Fenics quotes, two long and two short) the mean $\Cskew$ is 22.3 bp and the mean $\varphi$ is 0.64.

\subsection{E2: predictability of carry returns}\label{sec:e2}

\begin{table}[t]
\centering\footnotesize
\setlength{\tabcolsep}{4pt}
\caption{E2: predictive regressions of $\HML^U_{t+1}$}\label{tab:e2}
\begin{tabular}{lcccccc}
\toprule
\multicolumn{7}{l}{\emph{Panel A: in sample (159 months)}}\\
Predictor & $b$ & HAC $t$ & Bootstrap bias & Analytic bias & Corrected $b$ & One-sided $p$ \\
\midrule
$\varphi$ & 0.0062 & 5.95 & 0.00055 & 0.00049 & 0.0057 & 0.021 \\
Oriented $10\D$ RR & 0.397 & 4.13 & 0.0183 & 0.0148 & 0.379 & 0.004 \\
\midrule
\multicolumn{7}{l}{\emph{Panel B: out of sample (116 forecasts from January 2017)}}\\
Predictor & OOS $R^2$ & \multicolumn{2}{c}{Clark--West $t$ ($p$)} & \multicolumn{3}{c}{Clark--West $t$, Newey--West ($p$)} \\
\midrule
$\varphi$ & 3.7\% & \multicolumn{2}{c}{1.639 (0.051)} & \multicolumn{3}{c}{1.50 (0.067)} \\
Oriented $10\D$ RR & $-5.5$\% & \multicolumn{2}{c}{$-1.11$ (0.87)} & \multicolumn{3}{c}{$-1.18$ (0.88)} \\
\bottomrule
\end{tabular}
\par\medskip
\begin{minipage}{0.97\linewidth}\footnotesize
The dependent variable is the monthly return in decimal fractions; $\varphi$ is dimensionless and the risk reversal is in decimal volatility units. HAC $t$ is for the uncorrected $b$. Bootstrap bias from the residual bootstrap under the null (9,999 draws); analytic bias from Stambaugh (1999, eq.~18). The corrected $b$ subtracts the bootstrap bias; the one-sided $p$-value, for $b>0$, is from the same bootstrap. Out of sample: expanding window with training from the extended sample, one-sided test against the historical mean with the least-squares standard error; Newey--West alongside.
\end{minipage}
\end{table}

Table~\ref{tab:e2} reports E2. The two predictors are persistent (AR(1) coefficients of 0.86 for $\varphi$ and 0.52 for the risk reversal), with a correlation of 0.56, and their innovations are negatively correlated with returns ($-0.39$ and $-0.53$), so the least-squares slope is biased upwards in small samples (Stambaugh, 1999). After the bias correction, a higher skew price per unit of forward discount predicts higher carry returns in the following month (slope 0.0057, one-sided $p=0.021$), the sign a crash-compensation reading implies. The Newey--West $t$-statistic of 5.95 on the uncorrected slope is far stronger than the bootstrap evidence; for a persistent predictor with a heavy right tail it is unreliable, and it is not the pre-registered test.

Out of sample, the Clark--West statistic of 1.639 does not reach the one-sided 5\% critical value of 1.645 ($p=0.051$; 1.50 with the Newey--West standard error), so predictability is not established. Because training starts with the extended sample, the test splices the Fenics $\varphi$ of 2007--2013, built from two long and two short legs, with the composite primary series built from three and three; training from the primary start gives 0.96 (Section~\ref{sec:robust}). The risk reversal predicts in sample (corrected slope 0.379, $p=0.004$) but not out of sample.

The design also names option-implied variance and skewness as secondary predictors. Their moments are those of $\log(X_T/F_X)$ under the USD one-month forward measure, with dollar-base pairs converted by R1, from contracts spanned by out-of-the-money options (R3). Between the $10\D$ strikes the contracts come from the calibrated smile by composite Simpson's rule with the forward as a node, converged to $10^{-10}$ relative in all 1,440 currency-months; beyond them each contract is bounded by R4, and the variance and skewness intervals are the exact ranges over the box of contract intervals. Two tail settings were fixed before estimation: (i) the local elasticities of the smile's density at the $10\D$ strikes (median $\gamma=61$, $\eta=67$), which assumes the tails beyond the quotes are no heavier than at the quotes; and (ii) $\gamma=\eta=2$. The predictors are variance and skewness oriented to each leg (sign reversed for short legs), averaged over the six E1 legs.

Identification is weak. Under setting (ii) no skewness interval excludes zero, and the variance interval is about fifty times wider than under (i), so truncated quotes do not identify skewness without a tail assumption of the strength of (i). Under setting (i) the median currency-month skewness interval is $[-0.70,0.14]$, and 32\% of intervals exclude zero; but the calibrated smile's own extrapolation contradicts setting (i) in about three quarters of currency-months.

\begin{table}[t]
\centering\small
\caption{E2, secondary predictors: option-implied moments under tail setting (i)}\label{tab:mom}
\begin{tabular}{lccc}
\toprule
Predictor & $b$ (lower, mid, upper) & Newey--West $t$ & One-sided bootstrap $p$ \\
\midrule
Variance & 4.98, 4.81, 4.64 & 0.87, 0.90, 0.93 & 0.185, 0.171, 0.163 ($b>0$) \\
Oriented skewness & $-0.024$, $-0.028$, $-0.028$ & $-2.74$, $-2.60$, $-1.87$ & 0.013, 0.015, 0.045 ($b<0$) \\
\bottomrule
\end{tabular}
\par\medskip
\begin{minipage}{0.97\linewidth}\footnotesize
$\HML^U_{t+1}$, in decimal fractions, on the predictor evaluated at the lower endpoint, midpoint and upper endpoint of its R4 interval, 159 months. Variance is that of $\log(X_T/F_X)$; skewness is dimensionless. A conclusion is reported only if it holds at all three.
\end{minipage}
\end{table}

Table~\ref{tab:mom} reports the regressions under setting (i). Variance does not predict carry returns at any endpoint. More negative oriented skewness, meaning more left-tail risk in the carry position, predicts higher returns, the same economic direction as the test on $\varphi$. The one-sided bootstrap test rejects at 5\% at all three endpoints, but the Newey--West $t$ does not reach the two-sided 5\% critical value at the upper endpoint ($-1.87$). The result holds across the interval only under the bootstrap and only under a tail setting that the smile's own wings contradict, so it is conditional supporting evidence, not a finding.

\subsection{E3: decomposition of the realised hedge cost}\label{sec:e3}

\begin{table}[t]
\centering\small
\caption{E3: realised returns and the decomposition of the hedge cost}\label{tab:e3}
\begin{tabular}{lcc}
\toprule
Quantity & Mean & s.e. \\
\midrule
$\HML^U$ (unhedged) & 17.7 & 12.1 \\
$\HML^H$, $10\D$ & 8.3 & 11.9 \\
$\HML^H$, $25\D$ & 9.1 & 11.3 \\
$\HML^H$, at-the-money & 9.7 & 8.9 \\
\midrule
(i) payoff less $\sigP$ value & 0.5 & 3.8 \\
(ii) volatility level & 1.4 & 1.4 \\
(iii) skew (ex-ante premium) & $-11.2$ & 0.65 \\
\bottomrule
\end{tabular}
\par\medskip
\begin{minipage}{0.6\linewidth}\footnotesize
Primary sample, 159 months, bp per month, Newey--West standard errors. Terms (i)--(iii) are for the $10\D$ hedge.
\end{minipage}
\end{table}

Table~\ref{tab:e3} reports E3. Unhedged carry earns 17.7 bp per month with a standard error of 12.1, so the carry premium itself is not significant in this sample ($t=1.47$); the $10\D$ hedge lowers the mean to 8.3 bp. Term (iii), $-11.2$ bp per month, is the skew premium paid ex ante and carried to delivery, equal to the skew price of E1 by construction; its standard error of 0.65 measures the time variation of a premium known at the start of each month, not uncertainty about realised compensation, which is all in term (i). The point estimate of the realised hedge cost is close to the skew premium because the payoff term (0.5 bp, s.e.\ 3.8) and the volatility-level term (1.4 bp, s.e.\ 1.4) are indistinguishable from zero. The standard error of the total hedge cost, $\HML^H-\HML^U$, is not reported. By regime (supplementary, post hoc), unhedged carry averages 7.4 bp (s.e.\ 12.4) before 2022 and 37.0 bp (17.8) after, and hedged carry $-2.7$ bp (15.7) and 29.3 bp (16.0); neither difference is significant, with bootstrap 95\% intervals of $[-19.5,80.0]$ and $[-16.8,80.3]$ bp.

$\thUB$ is 0.53 at $10\D$, 0.49 at $25\D$ and 0.45 at the money, against a diffusive null of $\theta_0\approx0.10$ at $10\D$ (the average forward delta of the hedges) and about 0.25 for a $25\D$ put by R6. Its 95\% test-inversion set is unbounded, because the mean unhedged return is not significantly different from zero; as the design anticipated, the ratio is weakly identified, no value is claimed, and the gap between 0.53 and 0.10 is not read as a crash share. The realised terms also depend on the crashes that happened to occur in the sample, a peso-problem caveat: in the extended sample (secondary; Fenics quotes, 72 months, 2007--2013) term (i) is 16.7 bp (s.e.\ 10.2), reflecting option payoffs in 2008, the skew term is $-22.3$ bp (s.e.\ 3.0) and the unhedged mean is 8.4 bp with a standard error of 62.7.

\subsection{E5: systemic-risk exposure}\label{sec:e5}

\begin{table}[t]
\centering\small
\caption{E5: exposures to VIX roll-down and FX-volatility innovations}\label{tab:e5}
\begin{tabular}{lcccc}
\toprule
Portfolio & $\alpha$ (bp) & $\beta$, VIX roll-down & $\beta$, FX-volatility innovation & $R^2$ \\
\midrule
$\HML^U$ & $-7.3$ ($-0.61$) & 0.049 (7.41) & $-3.11$ ($-2.16$) & 0.33 \\
$\HML^H$, $10\D$ & $-10.6$ ($-0.79$) & 0.039 (4.69) & $-1.76$ ($-1.21$) & 0.22 \\
\bottomrule
\end{tabular}
\par\medskip
\begin{minipage}{0.9\linewidth}\footnotesize
Primary sample, 159 months. Newey--West $t$-statistics in parentheses. Returns are in decimal fractions and $\alpha$ in bp per month; the VIX roll-down is a decimal return and the FX-volatility innovation is in units of mean absolute daily log change. The AR(1) coefficient of the FX-volatility level is 0.76.
\end{minipage}
\end{table}

Table~\ref{tab:e5} reports E5. Unhedged carry returns load strongly on the VIX roll-down (correlation 0.55), consistent with Caballero and Doyle (2012), and negatively on FX-volatility innovations. The $10\D$ hedge lowers the VIX loading by about a fifth but leaves most of it, whereas Caballero and Doyle (2012) report that option-hedged carry retains high returns that systemic risk does not account for; the hedged FX-volatility loading is not significant. After adjusting for both factors, neither portfolio earns a significant $\alpha$; the exposure-adjusted ratio $1-\alpha_H/\alpha_U$ is $-0.45$ and its test-inversion set is unbounded, so it is unidentified. In the extended sample (secondary, 65 months) the loadings are larger, 0.124 unhedged ($t=5.39$) and 0.112 hedged ($t=4.53$), and the $\alpha$s are not significant.

\subsection{The regime-shift question under the pre-registered rules}

The design requires three conditions: an E1 difference whose interval excludes zero (not met: $-0.355$, bootstrap interval $[-0.814,0.014]$), a bias-corrected $b>0$ in E2 (met: 0.0057, one-sided $p=0.021$) and an out-of-sample Clark--West rejection at 5\% (not met: 1.639 against 1.645). Only the second holds, so the claim that the ex-ante price of crash insurance explains the change in the carry premium between the regimes has partial support only, the pre-registered label for this outcome: the condition met concerns in-sample predictability, and no estimate indicates that the price of crash insurance changed between the regimes. The skew-term rule of Section~\ref{sec:method} is uninformative here: term (iii) is an ex-ante premium, whose interval excludes zero whenever smiles are skewed, so its precision is mechanical and is not by itself evidence of priced crash risk.

\section{Robustness}\label{sec:robust}

\begin{table}[t]
\centering\footnotesize
\setlength{\tabcolsep}{2pt}
\caption{Robustness grid (one-month tenor)}\label{tab:rob}
\begin{tabular}{lcccccccc}
\toprule
 & \multicolumn{4}{c}{E1 and E2} & \multicolumn{4}{c}{Returns and E3 (bp per month)}\\
\cmidrule(lr){2-5}\cmidrule(lr){6-9}
Variant & $\D\varphi$ (s.e.) & Bootstrap 95\% & $b$ ($p$) & CW & $\HML^U$ & $\HML^H$ & Skew & $\thUB$ \\
\midrule
Base & $-0.355$ (0.204) & $[-0.82,0.01]$ & 0.0058 (0.017) & 0.96 & 17.7 (1.47) & 8.3 (0.70) & $-11.2$ & 0.53 \\
$25\D$ hedge & $-0.272$ (0.166) & $[-0.65,0.02]$ & 0.0074 (0.013) & 0.75 & 17.7 (1.47) & 9.1 (0.81) & $-8.4$ & 0.49 \\
ATM hedge & n.d. & & & & 17.7 (1.47) & 9.7 (1.09) & 0 & 0.45 \\
Smile strangle & $-0.345$ (0.196) & $[-0.80,0.01]$ & 0.0059 (0.018) & 0.99 & 17.7 (1.47) & 8.7 (0.73) & $-10.8$ & 0.51 \\
Fenics quotes & $-0.347$ (0.203) & $[-0.81,0.02]$ & 0.0058 (0.019) & 0.98 & 17.7 (1.47) & 8.5 (0.71) & $-11.2$ & 0.52 \\
Stale BF missing & $-0.440$ (0.265) & $[-1.04,0.02]$ & 0.0037 (0.023) & 0.19 & 17.7 (1.53) & 8.8 (0.76) & $-11.2$ & 0.50 \\
Excl.\ March 2020 & $-0.329$ (0.193) & $[-0.76,0.02]$ & 0.0053 (0.035) & 1.31 & 19.9 (1.73) & 8.8 (0.77) & $-11.0$ & 0.56 \\
Excl.\ CHF 2014--15 & $-0.353$ (0.204) & $[-0.82,0.01]$ & 0.0056 (0.016) & 0.84 & 20.4 (1.87) & 10.4 (0.89) & $-11.2$ & 0.49 \\
Implem., $k=1$ & base & base & 0.0059 (0.016) & 0.92 & 5.4 (0.45) & $-6.7$ ($-0.55$) & $-11.2$ & n.m. \\
Implem., $k=1.5$ & base & base & 0.0059 (0.016) & 0.92 & 5.4 (0.45) & $-8.0$ ($-0.66$) & $-11.2$ & n.m. \\
Implem., $k=2$ & base & base & 0.0059 (0.016) & 0.92 & 5.4 (0.45) & $-9.4$ ($-0.78$) & $-11.2$ & n.m. \\
Dollar carry & $-1.060$ (0.559) & $[-2.12,-0.37]$ & 0.0002 (0.373) & $-1.51$ & 0.3 (0.02) & $-2.6$ ($-0.17$) & $-2.8$ & n.m. \\
Ten currencies & $-0.300$ (0.149) & $[-0.63,-0.02]$ & 0.0053 (0.052) & 1.17 & 15.3 (1.37) & 8.2 (0.74) & $-9.7$ & 0.47 \\
Log returns & base & base & 0.0056 (0.020) & 0.86 & 16.7 (1.38) & 7.4 (0.62) & $-11.2$ & 0.56 \\
Previous-day spot & base & base & base & base & 17.7 (1.47) & 5.2 (0.43) & $-11.2$ & 0.70 \\
Vanna--volga & $-0.337$ (0.190) & $[-0.77,0.00]$ & 0.0060 (0.021) & 1.08 & 17.7 (1.47) & 8.7 (0.73) & $-10.8$ & 0.51 \\
\bottomrule
\end{tabular}
\par\medskip
\begin{minipage}{0.97\linewidth}\footnotesize
Primary sample; 160 month-ends for $\varphi$ and 159 return months, except the at-the-money (ATM) hedge (159), stale butterflies (BF) missing (156, 155) and excluding March 2020 (the February and March 2020 month-ends; 158, 157); CHF is excluded at the November 2014 to March 2015 month-ends. $\D\varphi$: hiking minus zero-rate, Newey--West s.e.\ in parentheses, stationary-bootstrap 95\% interval. $b$: bias-corrected E2 slope on $\varphi$, one-sided bootstrap $p$. CW: Clark--West statistic with training from the primary start (116 forecasts), not comparable with 1.639 in Table~\ref{tab:e2} (training from 2007). Returns: Newey--West $t$ in parentheses. Implem.: implementable, option half-spread $k$ times the ATM half-spread. base: as in the base row; n.d.: not defined; n.m.: not meaningful (mean unhedged return within half a standard error of zero). With 1,999 bootstrap draws, the base row's interval, $b$ and $p$ differ from Tables~\ref{tab:e1} and~\ref{tab:e2}.
\end{minipage}
\end{table}

\subsection{The one-month grid}

Table~\ref{tab:rob} reports the robustness grid. The rules the design leaves open were fixed in the estimation script after the primary results were recorded but before any variant was estimated. A currency is dropped at a month-end if its selected $25\D$ butterfly lies in a run of five or more identical two-sided business-day quotes (109 currency-months). A month-end is removed if its quote date or return window falls in an excluded period. Dollar carry follows Lustig, Roussanov and Verdelhan (2014, Section 3.2), long or short all available currencies with equal weights by the sign of their average forward discount (here the nine-currency average). In the ten-currency ranking the dollar enters with $\fd=0$ and a dollar leg has zero return and no option; the previous-day variant changes option payoffs only.

$\varphi$ is lower in the hiking regime in every variant. The difference is significant at 5\% under both inferences only in the ten-currency ranking ($t=-2.01$, interval $[-0.63,-0.02]$). For dollar carry only the bootstrap interval excludes zero, and its $\varphi$ divides by the absolute average forward discount, below 5 bp in 40 zero-rate months, so the difference comes mainly from the denominator. In the base and ten-currency portfolios the skew price per month is nearly the same in both regimes (11.2 and 11.2 bp; 10.0 and 9.3 bp), and the fall in $\varphi$ comes from the forward-discount spread. Under my reading of the rule the E1 conclusion therefore depends on the portfolio definition, although the ten-currency ranking is one of many correlated variants, with 1,999-draw intervals and no multiplicity adjustment. The smile-strangle reading and Fenics quotes move $\varphi$ by less than 4 per cent.

In E2, the bias-corrected slope is positive with one-sided $p\le0.035$ in every HML variant, $p=0.052$ for ten currencies and zero for dollar carry. No variant rejects out of sample at 5\% when training starts in 2013.

The skew term of E3 is $-8.4$ to $-11.2$ bp per month across the HML variants, and its small standard error reflects the stability of an ex-ante premium; $\thUB$ lies between 0.45 and 0.70 at mid prices, and the hedged mean is insignificant everywhere. Quoted spreads on forwards and spot reduce the unhedged mean from 17.7 to 5.4 bp, and option spreads make the hedged mean negative for every $k$. With the previous day's spot the hedged mean falls by 3.1 bp, from 43 months with an option in the money on one of the two days, no single month contributing more than 0.5 bp; neither close is the 10:00 New York cut, which lies between them, so both are reported.

\subsection{Vanna--volga smiles}

The alternative smile is the implied volatility of the vanna--volga price of Castagna and Mercurio (2007, eqs.~(6)--(7)), with Greeks at the at-the-money volatility; pillars at the delta-neutral-straddle strike and at the $25\D$ strikes at the pillar volatilities (Castagna and Mercurio, 2006, Section 2; Bossens, Ray\'ee, Skantzos and Deelstra, 2010, Sections 3.2 and 3.3); and, under the market reading, a smile strangle solved so that the smile prices the market strangle (Reiswich, 2010, Section 3.3.3; Bossens et al., 2010, Section 3.3). I chose the price smile before estimation, because the properties Castagna and Mercurio prove belong to the price; the first-order approximation overvalues the wings (Castagna and Mercurio, 2007, p.~10) and the second-order one can be undefined (Reiswich, 2010, Section 3.3.2). An implementation rule not from a source (a $10\D$ strike must connect to the at-the-money pillar through strikes where the smile is defined) does not bind. The smile passes the static-arbitrage check over $\pm4$ at-the-money standard deviations in 99.7\% of the 1,440 currency-months (market reading), and no leg uses one that fails. It gives $\varphi$ within 5 per cent of SABR in each regime, the same conclusions, and a better fit to the held-out $10\D$ risk reversal and butterfly (mean absolute errors of 0.10 and 0.13 volatility points, against 0.13 and 0.16 for SABR).

\begin{table}[t]
\centering\footnotesize
\setlength{\tabcolsep}{3pt}
\caption{Three-month tenor, quarterly non-overlapping positions}\label{tab:3m}
\begin{tabular}{lcccccccc}
\toprule
Phase & Quarters & $\D\varphi$ (s.e.) & Bootstrap 95\% & $b$ ($p$) & $\HML^U$ & $\HML^H$ & Skew & $\thUB$ \\
\midrule
Mar (primary) & 53; 35, 18 & $-0.301$ (0.166) & $[-0.63,\,-0.04]$ & 0.0228 (0.045) & 47.5 (1.29) & 36.6 (1.12) & $-24.9$ & 0.23 \\
Jan & 53; 34, 19 & $-0.275$ (0.167) & $[-0.63,\,-0.03]$ & 0.0142 (0.126) & 57.8 (1.64) & 42.8 (1.25) & $-23.6$ & 0.26 \\
Feb & 54; 35, 19 & $-0.246$ (0.130) & $[-0.52,\,-0.04]$ & 0.0119 (0.169) & 43.0 (1.18) & 9.7 (0.26) & $-22.7$ & 0.77 \\
\bottomrule
\end{tabular}
\par\medskip
\begin{minipage}{0.97\linewidth}
Phases by first quarter-end month: March, June, September, December (primary, from June 2013); January, April, July, October; February, May, August, November (both supplementary). Quarters of $\varphi$: total; zero-rate, hiking. $\D\varphi$ and $b$ as in Table~\ref{tab:rob}, with the Stambaugh bootstrap for $p$. Returns and skew term in bp per quarter, Newey--West $t$ in parentheses. Bootstraps use 1,999 draws.
\end{minipage}
\end{table}

\subsection{Three-month tenor}

Positions are quarterly and non-overlapping, three long and three short on the three-month forward discount, protected at the $10\D$ strike of the three-month smile. The primary rebalancing dates are the calendar-quarter month-ends from June 2013; the other two phases are supplementary. The quotes are complete for all nine currencies from May 2013 (11 of 11,520 fields substituted, 2 missing), and 1,436 of 1,440 smiles converge under each reading, all passing both arbitrage checks. In Table~\ref{tab:3m}, $\varphi$ falls from about 0.6 to 0.3 in every phase; the bootstrap interval excludes zero in all three, but the Newey--West $t$-statistic lies between $-1.65$ and $-1.89$, so under my reading the difference is not significant. The skew price per quarter changes little (25.5 to 23.6 bp in the primary phase) while the forward-discount spread rises from 55 to 82 bp. The bias-corrected E2 slope is positive in every phase, significant at 5\% only in the primary one ($p=0.045$); the skew term ($-23$ to $-25$ bp per quarter) is precisely estimated ($t$ between $-14$ and $-24$), and $\thUB$ (0.23 to 0.77) has an unbounded set in each phase.

\section{Conclusion}\label{sec:concl}

Under the pre-registered rules the answer is partial support only. The point estimate of $\varphi$ falls from 0.809 to 0.454, but the interval for the difference includes zero, and a supplementary (post hoc) split shows no evidence that the skew price per month changed (11.24 and 11.19 bp). $\varphi$ predicts carry returns in sample ($p=0.021$) but not out of sample (1.639 against 1.645): one condition of three, and it concerns predictability, not the change between regimes. The evidence points to a rise in the forward-discount spread rather than a change in the price of crash insurance.

The mean realised cost of $10\D$ protection is close to the skew premium paid ex ante, and unhedged carry is strongly exposed to the VIX roll-down, which the hedge lowers by about a fifth. The evidence cannot say what share of the carry premium compensates for crashes, or how the entropy difference behind it divides across cumulant orders (R7): the confidence sets of $\thUB$ (one and three months) and of the exposure-adjusted ratio are unbounded, because the unhedged carry premium is insignificant ($t=1.47$ at one month, 1.29 in the primary three-month phase), and quoted spreads reduce it from 17.7 to 5.4 bp. Under my reading of the E1 rule the conclusion depends on the portfolio definition, but in the ten-currency ranking, the only variant where both intervals exclude zero, the fall also comes from the forward-discount spread.

Several limitations qualify these conclusions. The hiking regime is short (56 month-ends), and a supplementary (post hoc) analysis finds roughly 8 effective observations per regime, so the pre-registered intervals are likely too narrow. The persistence and heavy right tail of $\varphi$ make the Newey--West inference in E2 unreliable, the out-of-sample result is sensitive to the training start, and the cross-section is small (nine currencies, three per side). $\thUB$ is weakly identified, and the realised E3 terms depend on the crashes that happened to occur. The $10\D$ premium is extrapolated from $25\D$ quotes, and SABR fits the held-out $10\D$ quotes less well than vanna--volga. Option-implied skewness is identified only under a tail setting that the smile's own wings contradict in about three quarters of currency-months. The quotes are dealer-contributed composites, not transaction prices, and implementable returns rest on an assumed spread rule. Payoffs use end-of-day spot, not the 10:00 New York cut; the hedged mean moves by 3.1 bp between the closes that bracket it. Dates use New York holidays only, no provider metadata confirm the description of the three-month USD OIS instrument, and quote revisions have not been checked.

Further data would narrow the remaining uncertainty. A longer hiking regime would add effective observations; executable single-option quotes would replace the spread rule; fixings at the 10:00 New York cut would settle payoff timing; quotes beyond $10\D$ would narrow the R4 intervals. Above all, a bounded set for $\thUB$ needs a significant unhedged carry premium, and so a longer sample of option-hedged carry.

\clearpage
\appendix

\section{Theory results}\label{app:theory}

Table~\ref{tab:theory} lists the results used in the paper. Proofs are in \texttt{theory/notes.tex}, with the hypotheses stated in full. Status: P, proved in the notes, with the part I claim as my own named in the last column; P$^*$, a known result whose published or standard proof the notes write out; C, cited with its hypotheses.

\begin{table}[H]
\centering\footnotesize
\caption{Theory results R1 to R9}\label{tab:theory}
\begin{tabularx}{\linewidth}{@{}l>{\raggedright\arraybackslash\hsize=1.25\hsize}X>{\raggedright\arraybackslash}p{1.9cm}>{\raggedright\arraybackslash\hsize=0.75\hsize}X@{}}
\toprule
 & Statement & Status & Source \\
\midrule
R1 & If $Z_t=S_tB^f_t/(S_0B^d_t)$ is a true martingale, $d\Q^f/d\Q^d=Z_T$ defines an equivalent foreign measure; Bayes formula; foreign and domestic martingale prices correspond; under Garman--Kohlhagen $\E^{\Q^d}S_T=F$ and $\E^{\Q^f}(1/S_T)=1/F$ & P$^*$ & Geman, El Karoui and Rochet (1995); Shreve (2004); strict local martingales: Carr, Fisher and Ruf (2014) \\
R2 & A function is the call-price curve of some $X\ge0$ with mean $F$ if and only if it is convex, equals $F$ at zero, has right slope at least $-1$ at zero and vanishes at infinity; the law is then unique & P$^*$ & F\"ollmer and Schied (2004, Lemma 7.23); context: Breeden and Litzenberger (1978); density test: Gatheral and Jacquier (2014) \\
R3 & Spanning of $\E f(X)$ by the forward and out-of-the-money puts and calls, under an explicit integrability condition & P$^*$ & Carr and Madan (2001), standard proof; contracts of Bakshi, Kapadia and Madan (2003), centred here at the forward \\
R4 & Bounds on the unobserved tails of put and call prices from observed wing prices and tail exponents; variance and skewness as intervals & P & Mine; truncation bounds for model-free variance: Jiang and Tian (2005); contracts adapted from Bakshi, Kapadia and Madan (2003) \\
R5 & (a) Closed-form strike for pips spot delta; (b) the premium-adjusted call delta rises to a unique maximum then falls, the convention taking the right root; (c) a slope condition makes the smile delta monotone, so inversion is unique & (a) standard; (b) P$^*$; (c) P & Reiswich and Wystup (2012) for definitions and the convention; (b): J\"ackel (2020), Mills-ratio monotonicity: Sampford (1953); (c): mine, with non-monotone smile deltas shown by Reiswich (2010) \\
R6 & Diffusive null of the hedge cost: $\theta_0=\Phi(-d_+(\xi))$ exactly for a forward $\xi$ between $F$ and $Fe^{\lambda\tau}$, with $d_+(G)=[\ln(G/K)+s^2/2]/s$ for the put's strike $K$; within $(2\pi)^{-1/2}|\lambda|\sqrt\tau/\sigma$ of $\Phi(-d_1)$, $d_1=d_+(F)$ & P (exact form and bound) & Leading-order form: Farhi et al.\ (2015); diffusive premium given up by an unlevered hedge: Jurek (2014) \\
R7 & Expected log currency excess return equals the entropy difference less the covered-interest-parity basis; cumulant expansion of entropy; entropy bound & P$^*$ & Backus, Foresi and Telmer (2001, eqs.~(9)--(13)); wedge: Lustig and Verdelhan (2019); convenience yields: Jiang, Krishnamurthy and Lustig (2021); Backus, Chernov and Zin (2014); bound: Bansal and Lehmann (1997) and Alvarez and Jermann (2005), as stated by Backus, Chernov and Zin; sign against Du, Tepper and Verdelhan (2018) \\
R8 & Hagan volatility with $\beta=1$ is smooth; the calibration has a locally unique $C^1$ solution in the quotes under non-degeneracy conditions & P & Hagan et al.\ (2002); market strangle: Reiswich and Wystup (2012); implicit-function-theorem argument: Reiswich (2010, Theorems 2--3); explicit leading-order inversion: Le Floc'h and Kennedy (2014); the local exact statement for SABR is mine \\
R9 & Central limit theorem for strongly mixing sequences (strict stationarity, $2+\delta$ moments, summable $\alpha(n)^{\delta/(2+\delta)}$, positive long-run variance); validity of the stationary bootstrap; bounded confidence sets for ratios have zero worst-case coverage, hence test inversion & C & Ibragimov (1962, Theorem 1.7); Politis and Romano (1994), Theorem 2 of the 1991 report; Gon\c{c}alves and de Jong (2003); Gleser and Hwang (1987); Dufour (1997); test inversion for a ratio: Fieller (1954) \\
GK & The Garman--Kohlhagen price solves its PDE and is the unique $C^{1,2}$ solution of polynomial growth; no boundary condition at $S=0$ is needed & P$^*$ & Garman and Kohlhagen (1983); It\^o and Feynman--Kac as in Shreve (2004) \\
\bottomrule
\end{tabularx}
\end{table}

\section{Reproduction and data access}\label{app:repro}

The code is in the repository; commands run from its root. The two Python environments are created with
{\small\begin{verbatim}
python -m venv .venv && .venv/bin/pip install -r requirements.lock \
  && .venv/bin/pip install -e . --no-deps
.venv/bin/python -m pytest -q
python -m venv .venv-lseg && .venv-lseg/bin/pip install -r requirements-lseg.lock
\end{verbatim}}
The LSEG acquisition reads the App Key of whoever runs it from \texttt{\string~/.lseg/app\_key}. The steps, in order, are those of the reports' reproduction blocks:
{\small\begin{verbatim}
# Acquisition and audit (Section 3; reports/data_audit.md)
.venv-lseg/bin/python scripts/acquire_lseg_fx.py --retrieval-date 2026-09-23
.venv/bin/python scripts/audit_fx_panel.py --retrieval-date 2026-09-23
# Smiles, E1 and E4 (Sections 4.1 and 5.1; reports/smile_calibration.md, e1.md)
.venv/bin/python scripts/calibrate_smiles.py --start 2010-07-30 --end 2026-08-31
.venv/bin/python scripts/calibrate_smiles.py --contributor FN \
  --start 2007-01-31 --end 2026-08-31
.venv/bin/python scripts/estimate_e1.py
# Returns, E2, E3, E5 and moments (Sections 5.2-5.4; reports/stage4.md)
.venv/bin/python scripts/estimate_stage4.py
.venv/bin/python scripts/acquire_cboe_vx.py --retrieval-date 2026-09-24
.venv/bin/python scripts/estimate_e5.py
.venv/bin/python scripts/estimate_moments.py
# Robustness (Section 6; reports/robustness.md)
.venv/bin/python scripts/calibrate_vv.py
.venv/bin/python scripts/robustness.py
.venv/bin/python scripts/calibrate_smiles.py --tenor 3M \
  --extra-raw-root data/private/lseg/2026-09-24/raw
.venv/bin/python scripts/estimate_3m.py
\end{verbatim}}
The three-month forward points and rates come from the separate retrieval of 24 September 2026, which the three-month calibration reads. The reports record no separate command for that retrieval; the acquisition script's instrument catalogue includes those instruments, and its \texttt{--only} and \texttt{--retrieval-date} options restrict a run to them and write it to that folder. The robustness script stops if its base variant does not reproduce the E1 regime difference and the means of $\HML^U$, $\HML^H$ and the skew term.

Rerunning the acquisition, and hence the estimation, requires the reader's own LSEG Workspace licence, not mine; the acquisition script lets a reader with that licence rebuild the private layer. Raw tables, cleaned data, audit counts by instrument, month-level series and calibrated parameters are kept in \texttt{data/private/}, which Git ignores, and a repository test fails if a data extract or anything under that folder is tracked. The Cboe settlement files are public but are stored privately because their redistribution terms are not established. Confirmation that the licence terms permit publication of pooled statistical summaries and coverage facts is pending.

\section*{References}
Page, section and equation numbers refer to the version consulted, as stated in each entry. ``Bibliographic record'' means that only the bibliographic details were checked and the source is cited for a standard result.

\smallskip
\begin{list}{}{\leftmargin=1.5em \itemindent=-1.5em \itemsep=2pt \parsep=0pt}\small
\item Alvarez, F. and Jermann, U. J. (2005). Using asset prices to measure the persistence of the marginal utility of wealth. \emph{Econometrica} 73(6), 1977--2016. Not consulted directly; cited as stated in Backus, Chernov and Zin (2014).
\item Backus, D., Chernov, M. and Zin, S. (2014). Sources of entropy in representative agent models. \emph{Journal of Finance} 69(1), 51--99. Version consulted: NBER Working Paper 17219, 2011.
\item Backus, D. K., Foresi, S. and Telmer, C. I. (2001). Affine term structure models and the forward premium anomaly. \emph{Journal of Finance} 56(1), 279--304. Version consulted: NBER Working Paper 5623, 1996.
\item Bakshi, G., Kapadia, N. and Madan, D. (2003). Stock return characteristics, skew laws, and the differential pricing of individual equity options. \emph{Review of Financial Studies} 16(1), 101--143.
\item Bansal, R. and Lehmann, B. N. (1997). Growth-optimal portfolio restrictions on asset pricing models. \emph{Macroeconomic Dynamics} 1(2), 333--354. Not consulted directly; cited as stated in Backus, Chernov and Zin (2014).
\item Bossens, F., Ray\'ee, G., Skantzos, N. S. and Deelstra, G. (2010). Vanna-volga methods applied to FX derivatives: from theory to market practice. \emph{International Journal of Theoretical and Applied Finance} 13(8), 1293--1324. Version consulted: arXiv:0904.1074v3, 3 May 2010.
\item Breeden, D. T. and Litzenberger, R. H. (1978). Prices of state-contingent claims implicit in option prices. \emph{Journal of Business} 51(4), 621--651. Consulted: bibliographic record.
\item Brunnermeier, M. K., Nagel, S. and Pedersen, L. H. (2008). Carry trades and currency crashes. \emph{NBER Macroeconomics Annual} 23, 313--348.
\item Burnside, C., Eichenbaum, M., Kleshchelski, I. and Rebelo, S. (2011). Do peso problems explain the returns to the carry trade? \emph{Review of Financial Studies} 24(3), 853--891.
\item Bussi\`ere, M., Chinn, M. D., Ferrara, L. and Heipertz, J. (2022). The new Fama puzzle. \emph{IMF Economic Review} 70(3), 451--486. Consulted: published abstract; authors' manuscript of 17 February 2022.
\item Caballero, R. J. and Doyle, J. B. (2012). Carry trade and systemic risk: why are FX options so cheap? NBER Working Paper 18644. Consulted: abstract, introduction and data section.
\item Carr, P., Fisher, T. and Ruf, J. (2014). On the hedging of options on exploding exchange rates. \emph{Finance and Stochastics} 18(1), 115--144. Version consulted: arXiv:1202.6188v3.
\item Carr, P. and Madan, D. (2001). Towards a theory of volatility trading. In E. Jouini, J. Cvitani\'c and M. Musiela (eds), \emph{Option Pricing, Interest Rates and Risk Management}, Cambridge University Press, 458--476. First published in R. Jarrow (ed.), \emph{Volatility}, Risk Books, 1998. Consulted: bibliographic record.
\item Castagna, A. and Mercurio, F. (2006). Consistent pricing of FX options. Working paper, Banca IMI; first posted as SSRN Working Paper 873788, January 2006. Version consulted: revision of 1 September 2006.
\item Castagna, A. and Mercurio, F. (2007). The vanna-volga method for implied volatilities. \emph{Risk}, January 2007, 106--111. Version consulted: authors' post-review preprint of 30 October 2006.
\item Cboe Futures Exchange. VX historical data: daily settlement files per contract. Retrieved 24 September 2026.
\item Chernov, M., Graveline, J. and Zviadadze, I. (2018). Crash risk in currency returns. \emph{Journal of Financial and Quantitative Analysis} 53(1), 137--170.
\item Choi, J. H. and Suh, S. (2022). Conditionally-hedged currency carry trades. \emph{Journal of International Financial Markets, Institutions and Money} 79, 101591. Consulted: published abstract and introduction only.
\item Clark, I. J. (2011). \emph{Foreign Exchange Option Pricing: A Practitioner's Guide}. Wiley. Not consulted directly; cited as reported in Reiswich and Wystup (2012).
\item Clark, T. E. and West, K. D. (2007). Approximately normal tests for equal predictive accuracy in nested models. \emph{Journal of Econometrics} 138(1), 291--311. Version consulted: NBER Technical Working Paper 326, 2006.
\item Della Corte, P., Ramadorai, T. and Sarno, L. (2016). Volatility risk premia and exchange rate predictability. \emph{Journal of Financial Economics} 120(1), 21--40. Version consulted: accepted manuscript, revised February 2015.
\item Du, W., Tepper, A. and Verdelhan, A. (2018). Deviations from covered interest rate parity. \emph{Journal of Finance} 73(3), 915--957. Version consulted: NBER Working Paper 23170, February 2017.
\item Dufour, J.-M. (1997). Some impossibility theorems in econometrics with applications to structural and dynamic models. \emph{Econometrica} 65(6), 1365--1387. Consulted: bibliographic record.
\item Fan, Z., Londono, J. M. and Xiao, X. (2022). Equity tail risk and currency risk premiums. \emph{Journal of Financial Economics} 143(1), 484--503.
\item Farhi, E., Fraiberger, S. P., Gabaix, X., Ranci\`ere, R. and Verdelhan, A. (2015). Crash risk in currency markets. Working paper, version of 12 March 2015, SSRN 1397668 (NYU Working Paper FIN-09-007); first issued as NBER Working Paper 15062, June 2009, revised May 2013. Versions consulted: SSRN version of 12 March 2015, to which section numbers refer, and NBER revision of May 2013.
\item Federal Reserve Board. Holidays Observed (K.8). \url{https://www.federalreserve.gov/aboutthefed/k8.htm}.
\item F\"ollmer, H. and Schied, A. (2004). \emph{Stochastic Finance: An Introduction in Discrete Time}, 2nd edition. De Gruyter. Consulted: Lemma 7.23 and its proof, in a page preview.
\item Fieller, E. C. (1954). Some problems in interval estimation. \emph{Journal of the Royal Statistical Society, Series B} 16(2), 175--185. Consulted: published article (JSTOR copy), Sections 1--4.
\item Garman, M. B. and Kohlhagen, S. W. (1983). Foreign currency option values. \emph{Journal of International Money and Finance} 2(3), 231--237. Consulted: bibliographic record.
\item Gatheral, J. and Jacquier, A. (2014). Arbitrage-free SVI volatility surfaces. \emph{Quantitative Finance} 14(1), 59--71. Version consulted: arXiv:1204.0646v4.
\item Geman, H., El Karoui, N. and Rochet, J.-C. (1995). Changes of num\'eraire, changes of probability measure and option pricing. \emph{Journal of Applied Probability} 32(2), 443--458. Consulted: bibliographic record.
\item Gleser, L. J. and Hwang, J. T. (1987). The nonexistence of $100(1-\alpha)\%$ confidence sets of finite expected diameter in errors-in-variables and related models. \emph{Annals of Statistics} 15(4), 1351--1362. Consulted: bibliographic record.
\item Gon\c{c}alves, S. and de Jong, R. (2003). Consistency of the stationary bootstrap under weak moment conditions. \emph{Economics Letters} 81(2), 273--278. Consulted: bibliographic record.
\item Hagan, P. S., Kumar, D., Lesniewski, A. S. and Woodward, D. E. (2002). Managing smile risk. \emph{Wilmott Magazine}, 84--108.
\item Ibragimov, I. A. (1962). Some limit theorems for stationary processes. \emph{Theory of Probability and Its Applications} 7(4), 349--382. Version consulted: Russian original, \emph{Teoriya Veroyatnostei i ee Primeneniya} 7(4).
\item J\"ackel, P. (2020). Strike from volatility and delta-with-premium. \emph{Quantitative Finance} 20(8), 1227--1235. Consulted: published article, Sections 1 and 2.
\item Jiang, G. J. and Tian, Y. S. (2005). The model-free implied volatility and its information content. \emph{Review of Financial Studies} 18(4), 1305--1342. Consulted: published article, Section 1.2.1.
\item Jiang, Z., Krishnamurthy, A. and Lustig, H. (2021). Foreign safe asset demand and the dollar exchange rate. \emph{Journal of Finance} 76(3), 1049--1089. Consulted: NBER Working Paper 24439, abstract and introduction.
\item Jurek, J. W. (2014). Crash-neutral currency carry trades. \emph{Journal of Financial Economics} 113(3), 325--347. Versions consulted: published abstract and introduction; SSRN 1262934, August 2013.
\item Kutuk, M. M. and van Wijnbergen, S. (2025). Carry trade and currency crash risk. CEPR Discussion Paper 20745. Version consulted: Tinbergen Institute Discussion Paper TI 2025-058/IV, 1 October 2025.
\item LSEG Developer Portal. IPA Volatility Surfaces: FX. Consulted 24 September 2026.
\item LSEG Workspace. FX option volatility, forward, deposit, OIS and VIX futures instruments. Retrieved 23 September 2026; three-month forwards and rates retrieved 24 September 2026.
\item Le Floc'h, F. and Kennedy, G. (2014). Explicit SABR calibration through simple expansions. Working paper, SSRN 2467231, version 1.2 of July 2014. Consulted: full text, Sections 1 and 4.
\item Lustig, H. and Verdelhan, A. (2019). Does incomplete spanning in international financial markets help to explain exchange rates? \emph{American Economic Review} 109(6), 2208--2244. Consulted: manuscript of January 2018, Section 2.5.
\item Lustig, H., Roussanov, N. and Verdelhan, A. (2011). Common risk factors in currency markets. \emph{Review of Financial Studies} 24(11), 3731--3777.
\item Lustig, H., Roussanov, N. and Verdelhan, A. (2014). Countercyclical currency risk premia. \emph{Journal of Financial Economics} 111(3), 527--553.
\item Menkhoff, L., Sarno, L., Schmeling, M. and Schrimpf, A. (2012). Carry trades and global foreign exchange volatility. \emph{Journal of Finance} 67(2), 681--718.
\item Newey, W. K. and West, K. D. (1987). A simple, positive semi-definite, heteroskedasticity and autocorrelation consistent covariance matrix. \emph{Econometrica} 55(3), 703--708. Consulted: bibliographic record.
\item Newey, W. K. and West, K. D. (1994). Automatic lag selection in covariance matrix estimation. \emph{Review of Economic Studies} 61(4), 631--653. Version consulted: NBER Technical Working Paper 144, 1993.
\item Patton, A., Politis, D. N. and White, H. (2009). Correction to ``Automatic block-length selection for the dependent bootstrap''. \emph{Econometric Reviews} 28(4), 372--375. Consulted: bibliographic record.
\item Politis, D. N. and Romano, J. P. (1994). The stationary bootstrap. \emph{Journal of the American Statistical Association} 89(428), 1303--1313. Version consulted: Technical Report 91-03, Department of Statistics, Purdue University, 1991.
\item Politis, D. N. and White, H. (2004). Automatic block-length selection for the dependent bootstrap. \emph{Econometric Reviews} 23(1), 53--70. Consulted: bibliographic record.
\item QuantLib, version 1.43 (Python bindings).
\item Reiswich, D. (2010). \emph{The Foreign Exchange Volatility Surface}. Doctoral dissertation, Frankfurt School of Finance \& Management. Version consulted: Deutsche Nationalbibliothek archive copy.
\item Reiswich, D. and Wystup, U. (2012). FX volatility smile construction. \emph{Wilmott} 2012(60), 58--69. Version consulted: CPQF Working Paper No.~20, version 2, Frankfurt School of Finance \& Management, 20 March 2010.
\item Sampford, M. R. (1953). Some inequalities on Mill's ratio and related functions. \emph{Annals of Mathematical Statistics} 24(1), 130--132. Consulted: bibliographic record.
\item Sheppard, K. (2026). arch. Python package, main branch at commit 91a25ef, 24 September 2026.
\item Shreve, S. E. (2004). \emph{Stochastic Calculus for Finance II: Continuous-Time Models}. Springer. Consulted: bibliographic record; cited for standard textbook statements.
\item Stambaugh, R. F. (1999). Predictive regressions. \emph{Journal of Financial Economics} 54(3), 375--421.
\item Zeileis, A. and contributors (2026). sandwich: Robust Covariance Matrix Estimators. R package version 3.1-3, 3 August 2026.
\end{list}

\end{document}
